%% kapa: exam=1 round=27 year=2016 period=2 form=A subject=appraisal_law subject_ko=감정평가관계법규 seq=22 number=22
%% source: Q-Net 감정평가사 기출문제 (cst00309) artlSeq=5169402 file=extracted/1008/5169402/제27회 감정평가사 1차 문제/제27회 감정평가사 2교시문제(A형).pdf
%% conversion: kordoc 4.24.51 -> tools/gen1.py (auto); review=proofread-sonnet; flags=none
\begin{question}{22}
국유재산법령상 행정재산에 관한 설명으로 옳은 것을 모두 고른 것은?
\begin{framedtext}
ㄱ. 행정재산의 사용허가를 받은 자가 그 행정재산의 관리를 소홀히 하여 재산상의 손해를 발생하게 한 경우에는 사용료 이외에 가산금을 징수할 수 있다. \par
ㄴ. 주거용으로 사용허가를 하는 경우에는 수의의 방법으로 사용허가를 받을 자를 결정할 수 없다. \par
ㄷ. 행정재산인 부동산에 대한 사용허가기간은 기부받은 재산을 그 기부자 등에게 사용허가하는 경우를 제외하고는 5년 이내로 한다. \par
ㄹ. 행정재산인 부동산을 직접 비영리 공익사업용으로 사용하려는 개인에게 사용허가한 경우 중앙관서의 장은 그 사용료를 면제할 수 있다. \par
\end{framedtext}
\begin{choices}
  \item ㄱ, ㄴ
  \item ㄱ, ㄷ
  \item ㄴ, ㄷ
  \item ㄴ, ㄹ
  \item ㄷ, ㄹ
\end{choices}
\end{question}
\answer{2}
